\FloatBarrier
\section{Related Work}
\label{sec:related_work}

\subsection{Model Collapse: Theory and Empirical Evidence}
\label{sec:rw_collapse}

Iterative training collapse was formalized by \citet{shumailov2024curse} and \citet{alemohammad2024mad}; \citet{schaeffer2025position} systematize eight distinct notions of collapse---our operationalization aligns with their \emph{distributional shift} family (definitions~D4--D6).
Theoretical analyses sharpen these results: \citet{dohmatob2024tale} recast collapse as changed scaling-law exponents, \citet{dohmatob2025strong} proved collapse at vanishingly small synthetic fractions, and \citet{truong2026phase} prove entropy collapse is a first-order transition, structurally ruling out critical-slowing-down early-warning signals~\citep{scheffer2009early,scheffer2012anticipating}.
Critical mixture ratios separating qualitative outcomes have independent precedent: \citet{kazdan2025collapse} map collapse conditions, \citet{gu2024cmr} establish scaling laws for critical ratios, and \citet{gu2025datamixing} show data mixing induces phase transitions in knowledge acquisition.

\subsection{Detection of Synthetic Data Contamination}
\label{sec:rw_detection}

Monitoring metrics proposed for collapse include expected calibration error (ECE)~\citep{du2025scales}, surplexity~\citep{gambetta2024surplexity}, and machine-generated text detection~\citep{drayson2025detection}; \citet{shi2025closer} independently validate entropy as a collapse diagnostic in diffusion models.
A complementary line of work pursues \emph{proactive} contamination prevention via watermarking~\citep{kirchenbauer2023watermark}; our distributional detection framework is \emph{reactive}---it identifies collapse from observed metric trajectories without requiring control over the generation process.
\citet{boettiger2012limits} establish fundamental power limits for early-warning detection in short time series ($T{<}25$), directly relevant to iterative training where observation windows are typically $T{=}5$--$20$.
However, these works evaluate single metrics in isolation without cross-validating against alternative statistical tests.
We differ by applying five detection methods to identical data, mapping detectability across contamination regimes, and treating method sensitivity as a primary finding rather than a robustness check (detailed comparison: Appendix~\ref{app:method_comparison}).

\subsection{Statistical Methods for Distribution Shift}
\label{sec:rw_methods}

For temporal precedence testing, Granger causality~\citep{granger1969investigating} is the standard framework, with neural extensions~\citep{tank2021neural,runge2019detecting}; however, non-stationary series inflate rejection rates~\citep{granger1974spurious}, and short series ($T{\leq}25$) preclude reliable nonlinear estimation~\citep{boettiger2012limits}.
Granger precedence establishes temporal predictive improvement, not true causation: bivariate tests cannot exclude confounders~\citep{shojaie2022granger}.
Our permutation-based approach circumvents the critical-slowing-down limitation by detecting temporal predictive improvement rather than variance or autocorrelation signatures.
Specification-curve analysis~\citep{simonsohn2020specification} formalizes the sensitivity of conclusions to analytical choices; we adapt this framework to show that method choice alone shifts detection rates from 0\% to 94\%.
The multiverse analysis perspective~\citep{steegen2016increasing,silberzahn2018many,coretta2023multidimensional,kapoor2024reforms} motivates our multi-method comparison: five methods on identical data produce conclusions ranging from universal significance to complete null, instantiating many-analysts divergence at the method level.
